% =====================================================================
%  PLANTILLA · presentación beamer (matemática) — suite LaTeX
%  Copiá este archivo y editá los frames.
%  Compilar:  latexmk -pdf beamer.tex
%  beamer YA carga hyperref, xcolor, amsmath y tikz (no recargarlos).
% =====================================================================
\documentclass[10pt,aspectratio=169]{beamer}

% --- Tema -----------------------------------------------------------
\usetheme{Madrid}
\usecolortheme{whale}
\setbeamertemplate{navigation symbols}{}   % oculta la barra de navegación

% --- Codificación y fuentes -----------------------------------------
\usepackage[utf8]{inputenc}
\usepackage[T1]{fontenc}
\usepackage{lmodern}
\usepackage{microtype}

% --- Idioma ---------------------------------------------------------
\usepackage[spanish,es-noshorthands]{babel}
\usepackage{csquotes}

% --- Matemática y contenido -----------------------------------------
\usepackage{mathtools}\usepackage{amssymb}\usepackage{bm}
\usepackage{siunitx}
\usepackage{booktabs}
\usepackage{pgfplots}\pgfplotsset{compat=1.18}

\hypersetup{unicode,   % marcadores/índice con acentos correctos
  pdftitle={Título de la presentación}, pdfauthor={Tu Nombre}}
\newcommand{\R}{\mathbb{R}}
\DeclareMathOperator{\sen}{sen}

% --- Metadatos del título -------------------------------------------
\title[Título corto]{Título de la presentación}
\subtitle{Subtítulo opcional}
\author{Tu Nombre}
\institute[UAI]{Universidad Abierta Interamericana}
\date{\today}

\begin{document}

\begin{frame}[plain]
  \titlepage
\end{frame}

\begin{frame}{Contenido}
  \tableofcontents
\end{frame}

% =====================================================================
\section{Primera sección}
% =====================================================================
\begin{frame}{Una idea por diapositiva}
  Texto principal del frame.
  \begin{itemize}
    \item Punto uno.
    \item Punto dos.
  \end{itemize}
\end{frame}

\begin{frame}{Bloques semánticos}
  \begin{block}{Definición}
    Un bloque normal para enunciar.
  \end{block}
  \begin{exampleblock}{Ejemplo}
    $f(x) = x^2$ es una parábola.
  \end{exampleblock}
  \begin{alertblock}{Atención}
    Un bloque de alerta para lo importante.
  \end{alertblock}
\end{frame}

\begin{frame}{Dos columnas}
  \begin{columns}[T]
    \begin{column}{0.48\textwidth}
      Columna izquierda: texto, ítems o ecuaciones.
      \[ \int_0^1 x^2 \, dx = \tfrac{1}{3}. \]
    \end{column}
    \begin{column}{0.48\textwidth}
      \begin{tikzpicture}
        \begin{axis}[width=\textwidth, height=4cm, axis lines=middle,
                     xlabel=$x$, samples=60]
          \addplot[blue, thick, domain=0:1] {x^2};
        \end{axis}
      \end{tikzpicture}
    \end{column}
  \end{columns}
\end{frame}

% Un frame con código verbatim/listings necesita [fragile]:
% \begin{frame}[fragile]{Con código}
%   \begin{verbatim}
%   codigo aqui
%   \end{verbatim}
% \end{frame}

\begin{frame}[plain]
  \centering \Large ¿Preguntas?
\end{frame}

\end{document}
